\par\noindent\textbf{(1)} 惯常每单位质量势写 Φ=−μ/r+[μJ₂R²/r³]P₂(sinϕ)，P₂(x)=(3x²−1)/2。轨道上 sinϕ=sin i sin(ω+f)。平均时利用 r=a(1−e²)/(1+e cosf) 和 dt∝r²df 得 ⟨Φ₂⟩=μJ₂R²[1−3cos²i]/[4a³(1−e²)\textasciicircum{}\{3/2\}]。
\par\noindent\textbf{(2)} 以 Delaunay 动量 L\_D=√(μa)、G\_D=L\_D√(1−e²)、H\_D=G\_D cosi 表示轴对称平均扰动，升交点 Ω\_node 的速率为 ∂⟨Φ₂⟩/∂H\_D（固定 G\_D）。求导得到 Ω̇\_node=−(3/2)nJ₂(R/p)²cosi，其中 n=√(μ/a³)、p=a(1−e²)；顺行轨道通常退行。
\par\noindent\textbf{(3)} 近心点角变量与 G\_D 共轭，求平均Hamiltonian对 G\_D 的偏导并注意 H\_D 固定、cosi=H\_D/G\_D，得到 ω̇=(3/4)nJ₂(R/p)²(5cos²i−1)。令其为零给临界倾角 i\_c=arccos(1/√5)≈63.435°（另有逆行对称角），可减少近心点长期漂移。
\par\noindent\textbf{(4)} 极轨 cosi=0 的节点漂移为零；顺行赤道形式有最大退行节点率但节点线本身不定义，应把极限与实体轨道区别；J₂→0 两个进动均消失。太阳同步轨道可用逆行倾角 cosi<0 产生正向 Ω̇\_node，从而匹配太阳视运动，而不仅靠调节高度。
